% Part: incompleteness
% Chapter: representability-in-q
% Section: prim-rec

\documentclass[../../../include/open-logic-section]{subfiles}

\begin{document}

\olfileid{inc}{req}{pri}
\olsection{Niru Rekursi Primitif}

Saiki kita bisa nuduhake yen definisi liwat rekursi primitif bisa
``ditiru'' nganggo panggolekan tanpa wates reguler liwat fungsi beta.
Upamana kita duwe $f(\vec x)$ lan $g(\vec x, y, z)$. Mula fungsi~$h(x,\vec
z)$ kang ditegesi saka $f$ lan~$g$ liwat rekursi primitif yaiku
\begin{align*}
h(\vec x, 0) & =  f(\vec x) \\
h(\vec x, y+1) & =  g(\vec x, y, h(\vec x, y)).
\end{align*}
Kita kudu nuduhake yen $h$ bisa ditegesi saka $f$ lan~$g$ mung nganggo
komposisi lan panggolekan tanpa wates reguler, kanthi nggunakake
fungsi dhasar lan fungsi kang ditegesi saka fungsi mau nganggo
komposisi lan panggolekan tanpa wates reguler (kayata~$\beta$).

\begin{lem}
\ollabel{lem:prim-rec}
Yen $h$ bisa ditegesi saka $f$ lan $g$ nganggo rekursi primitif, iku
bisa ditegesi saka $f$, $g$, lan fungsi $\Zero$, $\Succ$,
$\Proj{n}{i}$, $\Add$, $\Mult$, $\Char{=}$, nganggo komposisi lan
panggolekan tanpa wates reguler.
\end{lem}

\begin{proof}
Luwih dhisik, tegesna fungsi bantu $\hat h(\vec x, y)$ kang mbalekake
bilangan paling cilik~$d$ saengga $d$ ngodeake urutan kang nyukupi
\begin{enumerate}
\item $(d)_0 = f(\vec x)$, lan
\item kanggo saben $i < y$, $(d)_{i+1} = g(\vec x, i, (d)_i)$,
\end{enumerate}
ing kene $(d)_i$ iku cekakan saka $\beta(d,i)$. Kanthi tembung liya, $\hat h$
mbalekake urutan $\tuple{h(\vec x, 0), h(\vec x, 1), \dots, h(\vec
x, y)}$. Kita bisa nulis $\hat h$ dadi
\[
\hat h(\vec x, y) = \umin{d}{(\beta(d,0) = f(\vec x) \land \bforall{i <
  y}{\beta(d,i+1) = g(\vec x, i,\beta(d,i)})}.
\]
Gatekna: ing kene ora perlu rekursi primitif, mung minimalisasi. Fungsi
kang diminimalisasi iku reguler amarga lema fungsi beta
\olref[bet]{lem:beta}.

Nanging saiki kita nduweni
\[
h(\vec x, y) = \beta(\hat h(\vec x, y), y),
\]
mula $h$ bisa ditegesi saka fungsi dhasar mung nganggo komposisi
lan panggolekan tanpa wates reguler.
\end{proof}

\end{document}
